\section{The Feedback Loop Survival Guide}
\textit{Reference $|$ Course Text: Chapter 13}

Feedback Loop peer evaluations directly affect your individual grade. SDI has 2 evaluations; SDII has 3.

\begin{tipbox}
There are two types of individual assessment: \textbf{Peer Evaluations} (completed by your teammates about each other) and \textbf{Individual Performance Evaluations} (completed by your PA about you). Both inform your individual grade. Check Canvas for the specific schedule of each.
\end{tipbox}

\begin{warnbox}
A pattern of low Feedback Loop scores (below 3.0 out of 5.0) combined with PA observation of non-participation can result in individual grade reduction, PIP initiation, or course failure regardless of team scores.
\end{warnbox}

\textbf{How to receive good evaluations}: attend every meeting on time, complete your tasks by the agreed deadline, communicate proactively when you are behind, help teammates when you finish early, participate actively in discussions (not just silently), and be responsive to messages within 24 hours.

\textbf{How to give useful evaluations}: be specific. ``Jordan was great'' does not help the PA calibrate. ``Jordan completed the firmware ahead of schedule, debugged the I2C communication issue that was blocking the sensor team, and volunteered to write the CDR power analysis section when the original author was overwhelmed'' gives the PA evidence to work with.

\textbf{Self-evaluation}: be honest. PAs cross-check your self-rating against peer ratings and their own observations. Over-rating yourself while peers rate you lower signals a lack of self-awareness. Acknowledging a specific weakness (``I was late on two deliverables in Sprint~10 due to competing course deadlines; I should have communicated the conflict earlier'') signals maturity.


% ═══════════════════════════════════════
